% !TEX root = ../macro.tex
\section{Fluctuations: Mechanisms, Comparisons, and Diagnosis}
\label{sec:fluctuations-synthesis}

A district reports fewer orders, shorter shifts, falling sales, and postponed
investment. Has its ability to produce declined, have customers postponed
purchases, or can otherwise viable plans no longer be financed? Those questions
can lead to different models even when the headline observation is the same.
A model becomes useful when it identifies the constraint that changed and the
additional observations that would distinguish its explanation.

The district in this chapter is hypothetical. We follow its firms and workers
through several possible situations rather than assign one model to a country
or a recession. The purpose is to organize what has been derived, show how the
mechanisms fit together, and identify where a new model is required. The
financial and matching mechanisms discussed at the boundary are questions for
later development, not results already established by the completed chapters.

\subsection{One framework, several questions}

Growth models were easier to separate because their common question was how
income and productive capacity develop over time. Saving, intertemporal choice,
generations, and technical change supplied distinct routes into that question.
For fluctuations, the same optimizing household can face a technology change,
a change in the real return on saving, an inherited contract, or a borrowing
limit. Naming the model is only the beginning of identifying the mechanism.

Keep four questions in view. \emph{Source}: what changes first? \emph{Propagation}:
which stocks, contracts, or expectations carry its effects forward?
\emph{Adjustment}: which prices or quantities can move, and which cannot?
\emph{Closure}: which rule and market conditions make all the plans consistent?
The shock, its persistence, the constraint, and the policy response are separate
parts of the explanation.

\begin{center}
\begin{minipage}{\linewidth}
\captionsetup{type=table}
\centering\small
\begin{threeparttable}
\caption{The questions isolated by the completed fluctuation models.}
\label{tab:fluctuation-mechanism-map}
\begin{tabular}{p{.18\linewidth}p{.35\linewidth}p{.37\linewidth}}
\hline
\tableheading{Specification}&\tableheading{Question made visible}&\tableheading{Adjustment and memory}\\
\hline
RBC&How do households and firms reallocate work, spending, and investment when real opportunities change?&Flexible real prices coordinate choices. Capital carries history in the accumulation version; shock persistence and forecasts also matter.\\
Sticky-price NK&Why can a change in spending or monetary policy move output relative to flexible-price output?&Some product quotes are inherited. Resetters consider future costs. The basic consumption-only gap system has no inherited capital.\\
Sticky-wage NK&What changes when compensation cannot always respond to new work incentives and labor costs?&Wage quotes are inherited as well. The real-wage position connects price and wage inflation and adds an endogenous state.\\
\hline
\end{tabular}
\end{threeparttable}
\end{minipage}
\end{center}

Table~\ref{tab:fluctuation-mechanism-map} contains specifications within a common
dynamic general-equilibrium approach, not three unrelated languages. The wage
chapter already combines sticky prices with sticky wages. DSGE describes the
framework; NK describes a set of economic mechanisms within it. Adding capital,
financial intermediaries, or heterogeneous households expands the specification
and requires additional budgets and choices. It does not turn a label into a
solution by itself.

\begin{greenblock}
\textbf{Remember the constraint before the model name.}
Ask who wants to change a decision, what prevents that change, and which other
variable then adjusts. A price contract, a wage contract, an inherited machine,
and a borrowing limit are different constraints. Their effects depend on the
shock and the policy environment in which they operate.
\end{greenblock}

\subsection{The same output decline can describe different economies}

\subsubsection{When production opportunities deteriorate}

First suppose the district loses access to a productive input. A machine runs
more slowly or a production process becomes less effective. In the RBC language,
the productivity state changes. The basic production accounting is
\begin{align}
 Y_t&=A_tK_t^{\alpha}N_t^{1-\alpha},\\
 \hat y_t&=a_t+\alpha\hat k_t+(1-\alpha)\hat n_t.
\end{align}
At impact, the inherited capital stock cannot jump. If it was initially at
reference, its deviation is zero, so
\begin{align}
 \hat y_0&=a_0+(1-\alpha)\hat n_0.
\end{align}
Output changes through productivity and the endogenous hours response. Lower
productivity does not mechanically determine every other sign: income effects,
work substitution, expected recovery, and the investment opportunity all enter
the household's decision. The solved RBC experiment supplies a particular
joint response under its stated parameters.

Before treating the decline as spare capacity, ask whether the district can
still produce its previous quantity with the same resources. If not, restoring
spending alone does not restore the missing production opportunity. Conversely,
a temporary decline in orders does not prove that the production function
shifted. The resource constraint and the demand schedule answer different
questions.

RBC's real allocation also does not, by itself, supply a prediction for the
nominal price level or inflation. To interpret those observations, a nominal
closure must be specified. Describing a fall in technology as ``output down,
inflation up'' without such a closure adds a result that the real model has
not produced.

\subsubsection{When customers postpone purchases}

Now leave the district's productive capacity unchanged. Households want to
postpone consumption because the real return is higher relative to the
natural rate. Gap IS says
\begin{align}
 x_t&=E_tx_{t+1}-\frac1\sigma(r_t-r_t^n),\qquad
 r_t=j_t-E_t\pi_{t+1}^p.
\end{align}
At a held next-gap forecast, a higher real-rate gap lowers current demand.
The NK equilibrium also determines how the forecast, inflation, and policy
respond; the conditional equation alone is not the entire shock solution.

A restaurant with inherited menu prices can meet fewer bookings by reducing
meals and shifts. Marginal cost and new price quotes then respond. This is the
mechanism isolated by the policy experiment, not evidence that every restaurant
with weak sales faces a monetary disturbance. A preference change could alter
the natural rate, and expected future income could change spending independently
of a current policy innovation.

The question is whether the spending path is weak relative to the model's
flexible-price allocation. A measured trend gap is not automatically that
counterfactual. Technology and preferences can move the natural allocation
while observed output changes as well.

\subsubsection{When labor compensation adjusts on a different schedule}

Suppose orders fall before annual pay reviews. Wages may remain close to their
old nominal quotes while work incentives and product prices change. The relevant
comparisons are
\begin{align}
 \hat{mc}_t&=g_t+\gamma_p x_t,\\
 \hat\mu_t^w&=g_t-\gamma_wx_t.
\end{align}
A real wage that is only slightly below its natural reference can still be
high relative to the MRS associated with lower consumption and fewer hours.
Table~\ref{tab:sw-impact-comparison} illustrates exactly that combination:
the both-contract experiment has a small negative real-wage gap and a positive
wage-markup deviation.

This distinction is useful when nominal pay changes little but real purchasing
power falls. It is also why the response cannot be diagnosed from a real-wage
series alone. The consumer price, hours, household valuation of consumption,
and the compensation schedule all enter the comparison.

\subsubsection{When the missing constraint is outside these models}

A supplier may have customers and working machines yet be unable to finance
materials. A worker may be willing to work at the going wage yet fail to find
a matching vacancy. Neither situation is fully described by the representative
household's interior hours choice. A financing constraint requires assets,
liabilities, repayment, and a rule limiting funds. Job finding requires workers,
vacancies, and a matching or recruitment process.

The financial and labor-market chapters are where those conditions belong.
Here the observations tell us which missing mechanism to seek. Renaming a
monetary innovation ``credit stress'' does not introduce a lender's balance
sheet; renaming lower hours ``unemployment'' does not create job search.

\begin{center}
\begin{minipage}{\linewidth}
\captionsetup{type=table}
\centering\small
\begin{threeparttable}
\caption{Further observations that help distinguish explanations of weaker activity.}
\label{tab:fluctuation-diagnosis}
\begin{tabular}{p{.19\linewidth}p{.33\linewidth}p{.38\linewidth}}
\hline
\tableheading{Candidate constraint}&\tableheading{What to look at next}&\tableheading{What the observation would still not establish}\\
\hline
Production opportunity&Input availability, output per hour, usable capital, and production interruptions&A productivity residual can include utilization or measurement changes; it is not a direct count of technical ideas.\\
Desired spending&Orders, inventories, spending plans, expected income, and expected real returns&Weak sales alone do not identify a monetary shock or hold supply conditions fixed.\\
Nominal contracts&Timing of price and wage reviews, relative quotes, real wages, and hours&An unchanged wage does not imply an unchanged real wage or a particular unemployment mechanism.\\
Financing&Collateral, spreads, lending conditions, repayment needs, and rejected projects&A wider spread can reflect risk or demand as well as funding constraints. An explicit intermediary model is still needed.\\
Matching and reallocation&Vacancies, job finding, separations, skills, and participation&Slow hiring can reflect several constraints; wage stickiness alone does not measure them.\\
\hline
\end{tabular}
\end{threeparttable}
\end{minipage}
\end{center}

These are hypotheses for a diagnostic sequence, not sign restrictions that
identify a unique cause from one observation. Several constraints can bind in
the same district and at different stages of a downturn.

\subsection{A short-lived disturbance can leave a longer history}

A useful memory from Growth is the difference between an inherited stock and
a current decision. Capital survives today's saving decision. In the wage
model, an inherited relative compensation position survives today's policy
innovation. Neither persistence needs to be assigned entirely to the shock.

For the monetary experiment in the wage chapter,
\begin{align}
 \nu_t&=\rho_\nu\nu_{t-1}+\varepsilon_t^\nu,\\
 g_t&=Gg_{t-1}+J\nu_t,\qquad
 x_t=f_xg_{t-1}+h_x\nu_t.
\end{align}
Apply one date-zero innovation $e$ and zero later innovations, with zero
prehistory. Then $\nu_t=\rho_\nu^te$. Iterating the state law gives
\begin{align}
 g_0&=Je,\\
 g_1&=GJe+J\rho_\nu e,\\
 g_t&=Je\sum_{h=0}^tG^{t-h}\rho_\nu^h.
\end{align}
To evaluate this sum for $G\ne\rho_\nu$, multiply it by $G$ and by
$\rho_\nu$, then subtract:
\begin{align}
 G\sum_{h=0}^tG^{t-h}\rho_\nu^h
 -\rho_\nu\sum_{h=0}^tG^{t-h}\rho_\nu^h
 &=G^{t+1}-\rho_\nu^{t+1},\\
 g_t&=Je\frac{G^{t+1}-\rho_\nu^{t+1}}{G-\rho_\nu}.
\end{align}
If $G=\rho_\nu$, every term of the original sum is $G^t$, so instead
\begin{align}
 g_t&=Je(t+1)G^t.
\end{align}
The two roots describe different memories: $\rho_\nu$ belongs to the policy
process and $G$ to the endogenous contract adjustment. In the iid case,
$\rho_\nu=0$, the disturbance vanishes after impact but
\begin{align}
 g_t&=JeG^t,\qquad
 x_t=f_xJeG^{t-1}\quad(t\geq1).
\end{align}
The later response need not have the same sign as the impact response.
Persistent adjustment and permanently lower output are also different claims.
With stable roots, this model's gap eventually vanishes. A destroyed capital
stock, a different trend, or hysteresis requires its own state equation.

\begin{blueblock}
\textbf{Three different questions about persistence.}
How long does the exogenous disturbance last? How long does the economy take
to adjust after it ends? Does the eventual level or growth path change?
A long recovery does not, by itself, answer the third question. Nor does a
short-lived innovation imply a short-lived equilibrium response.
\end{blueblock}

\subsection{Compare the same policy disturbance under different constraints}

Figure~\ref{fig:sw-contract-responses} holds the policy innovation fixed and
solves each contract specification separately. The differences therefore reflect
changed adjustment mechanisms and the resulting forecasts, not independently
chosen response curves. Its output comparison concerns the same labor-only
family of models; it is not a comparison between a full capital RBC calibration
and an unrelated estimated NK economy.

A flexible-price-and-wage limit clarifies what the contracts add. With no
real shocks and constant markup wedges, flexible quotes give $x_t=g_t=0$.
Dynamic IS therefore requires
\begin{align}
 0&=0-\frac1\sigma(j_t-E_t\pi_{t+1}^p),\\
 j_t&=E_t\pi_{t+1}^p.
\end{align}
Combine that condition with the same Taylor rule:
\begin{align}
 \phi_\pi\pi_t^p+\nu_t&=E_t\pi_{t+1}^p.
\end{align}
For $\nu_t$ following AR(1), try $\pi_t^p=b_\nu\nu_t$:
\begin{align}
 \phi_\pi b_\nu\nu_t+\nu_t&=b_\nu\rho_\nu\nu_t,\\
 (\phi_\pi-\rho_\nu)b_\nu&=-1,\\
 \pi_t^p&=-\frac{\nu_t}{\phi_\pi-\rho_\nu},\qquad
 j_t=-\frac{\rho_\nu\nu_t}{\phi_\pi-\rho_\nu}.
\end{align}
For $\phi_\pi>1$ and $0\leq\rho_\nu<1$, this is the bounded solution.
The nominal variables adjust while the real allocation stays at its natural
benchmark. A positive \emph{rule intercept} need not produce a higher realized
nominal rate when inflation responds. This makes it especially important to
separate a policy innovation, the systematic reaction, and the observed rate.

With a sticky product or wage contract, the same nominal innovation can also
move the real-rate gap and current output. The sign and persistence then follow
from the complete system, including expectations and the inherited states.
The conditional IS--LM diagrams help organize those channels; the solved dynamic
paths determine how their intercepts move over time.

\subsection{Read a recession as a sequence of constraints}

Return to the district and begin with a sequence rather than a model label.
Did a supplier lose an input before it lost orders? Did orders weaken before
staffing changed? Were prices reviewed before pay agreements? Did customers
cancel because they wanted to postpone purchases, or because their lender
withdrew funds? The ordering does not identify causality by itself, but it
helps formulate a model in which the proposed cause can actually operate.

Next follow the accounts. For households, a wage reduction is also an income
change; for firms it is a cost change. A loan is another sector's asset. Lower
investment reduces current demand and changes future capital. A model that
keeps only one side of these transactions may describe a partial mechanism
while missing the aggregate feedback. General equilibrium makes those sides
fit together; it does not guarantee that their combined effect is desirable.

Then separate the horizon of the question. A temporary inability to sell meals
can leave the restaurant's productive knowledge intact. A prolonged inability
to finance maintenance can affect later capacity. A worker missing a shift and
a worker losing a durable employment match carry different histories. The
relevant horizon follows from what is inherited and what is costly to rebuild,
not from an arbitrary rule that fluctuations must disappear after a fixed
number of quarters.

Finally ask for an observation that could make the explanation fail. If the
story requires unchanged technology, what evidence of input loss would challenge
it? If it requires an inherited wage, what would immediate repricing of labor
change? If it requires binding credit, would a sufficiently funded firm in the
same market respond differently? This is how an intuition becomes a model
comparison rather than a flexible narrative that fits every outcome.

\begin{redblock}
\textbf{A disciplined comparison.}
Choose the initial state, the disturbance, the variables observed, and the
policy rule. Change one adjustment mechanism, re-solve the equilibrium, and
explain why the resulting joint paths differ. If instead the shock or policy
also changes, say so explicitly. A model-generated response is a conditional
prediction; identifying the disturbance in data is a further task.
\end{redblock}

\subsection{A smaller gap is not the whole welfare question}

The natural allocation allows flexible quotes while retaining market power
and other specified real distortions. The efficient allocation removes the
relevant wedges under a stated planner problem. A statistical trend is a fitted
reference path. Productive capacity describes feasible production under a
chosen resource concept. These objects need not coincide, and no one output
series supplies all of them.

The earlier chapters showed that constant price and wage markups depress
natural output relative to the symmetric efficient allocation. Contract dispersion
also changes how spending and hours are allocated across varieties and labor
types. Stabilizing an aggregate gap alone therefore need not remove the
welfare cost of unequal relative prices and hours.

A small calculation illustrates one work-allocation cost. Take two equally
weighted labor types with physical hours $H(1+z)$ and $H(1-z)$, so average
physical hours stay $H$. This is an illustration of the work burden at a held
arithmetic mean, not an equilibrium with composite output also held fixed.
The average disutility is
\begin{align}
 \mathcal V(z)&=\frac{\chi H^{1+\varphi}}{2(1+\varphi)}
                       [(1+z)^{1+\varphi}+(1-z)^{1+\varphi}].
\end{align}
Expand the two powers separately:
\begin{align}
 (1+z)^{1+\varphi}
 &\simeq1+(1+\varphi)z+\tfrac12\varphi(1+\varphi)z^2,\\
 (1-z)^{1+\varphi}
 &\simeq1-(1+\varphi)z+\tfrac12\varphi(1+\varphi)z^2,\\
 \mathcal V(z)
 &\simeq\frac{\chi H^{1+\varphi}}{1+\varphi}
                       +\frac{\chi\varphi H^{1+\varphi}}2 z^2.
\end{align}
The first-order terms cancel, yet uneven work raises the burden at second
order when $\varphi>0$. This explains why a first-order dispersion cancellation
does not justify omitting dispersion from a welfare analysis. The full policy
loss function additionally needs the output allocation, product-price dispersion,
reference distortions, and explicit instruments.

\subsection{How the remaining parts extend the map}

The planned labor-market chapter distinguishes hours from employment, job
finding, vacancies, and participation. It adds a mechanism that the completed
Calvo-wage model does not contain. Open Economy then allows the district to
trade goods and future income with others: demand and financing can move
across borders, relative prices include exchange rates, and adjustment depends
on the monetary arrangement.

The money, finance, and policy part adds balance sheets, borrowing constraints,
and household differences before comparing instruments and objectives. These
are also fluctuation mechanisms even though the book develops them later.
The part boundary organizes the learning sequence; it is not a claim that
finance or international trade matters only outside a recession model.

\begin{greenblock}
\textbf{A map to carry into the next part.}
Describe the observation, identify the candidate disturbance, name the constraint,
trace the inherited state, and specify policy. Then compare the implied paths
with other mechanisms and with evidence. This sequence links the growth models,
the real and nominal fluctuation benchmarks, and the extensions still to come.
\end{greenblock}
